\chapter{Failure Modes: как память ломается}
\label{ch:failuremodes}

\section{Зачем нужен отдельный каталог поломок}

У сложной памяти не бывает одного режима отказа. Она может ошибаться в содержании,
в топологии, в темпоральной динамике, в provenance, в контекстной упаковке и в
модальности ответа. Если всё это смешать, любая ошибка выглядит одинаково:
``система ответила странно''. Отдельный каталог нужен, чтобы странность стала
диагностируемой.

\section{Топологические поломки}

\subsection*{Hub collapse}
Слишком высокая концентрация $\pi$ на хабах приводит к тому, что память отвечает
через общие, часто посещаемые узлы, а не через редкие релевантные.

\textbf{Симптомы}: однообразные мосты, частое использование общих концептов,
потеря частного контекста.

\textbf{Наблюдаемые переменные}: корреляция между частотой включения узла в ответ
и его степенью; распределение visits.

\textbf{Предохранители}: штраф за сверхцентральные промежуточные узлы, рестарт на
якорях, отдельный glue budget.

\subsection*{Bottleneck starvation}
Нужные контрсвидетельства сидят за узким сечением графа и не попадают в retrieval.

\textbf{Симптомы}: система упорно повторяет прежнюю картину мира, хотя в графе уже
есть коррекция.

\textbf{Предохранители}: телепорт-скраб, явные bridge-edge policies, targeted
forensic retrieval.

\section{Эпистемические поломки}

\subsection*{Premature certainty}
Гипотеза слишком рано превращается в ``почти факт''.

\textbf{Симптомы}: мало узлов в зоне незнания, низкая доля осторожных формулировок,
частая болезненная ретракция.

\textbf{Предохранители}: priors по типам выводов, penalty за зависимые источники,
запрет на $|\conf|=1$ у выводных узлов.

\subsection*{Provenance erosion}
После нескольких циклов консолидации память ещё хранит утверждения, но уже не может
объяснить, откуда они взялись.

\textbf{Симптомы}: ответы звучат уверенно, но support-skeleton отсутствует.

\textbf{Предохранители}: отдельная метрика provenance coverage, обязательная
передача support-skeleton при prune и summarize.

\section{Динамические поломки}

\subsection*{Mood lock-in}
Настроение слишком сильно переопределяет маршруты retrieval, и память начинает
самоусиливать текущее состояние.

\textbf{Симптомы}: контекст ответа систематически смещён в одну валентность,
возникают руминационные петли.

\textbf{Предохранители}: saturation на mood-congruence, аварийный reflective mode,
ограничение числа self-facts, требование независимого подтверждения.

\subsection*{Over-consolidation}
Память слишком активно сворачивает частные эпизоды в обобщения и теряет материал
для будущей ревизии.

\textbf{Симптомы}: красивые summaries, но пустые forensic trails; исчезновение
конкретных различающих эпизодов.

\textbf{Предохранители}: хранение support-skeleton, delayed prune, тесты на
ability-to-disconfirm.

\section{Инженерные поломки}

\subsection*{Context bloat}
Retrieval находит слишком много всего, а context packer не умеет экономить.

\textbf{Симптомы}: длинный, тяжёлый, но рыхлый контекст; LLM тратит слот на чтение
мусора вместо решения.

\textbf{Предохранители}: dual budget packer, quota на provenance, лимиты по типам узлов.

\subsection*{Pending avalanche}
В thought-inbox копится больше кандидатов, чем коммиттер может осмысленно судить.

\textbf{Симптомы}: растущий хвост необработанных гипотез, устаревание pending-материала,
деградация explainability.

\textbf{Предохранители}: aging pending-кандидатов, priority queues, дешёвый triage,
quota на генерацию новых fusion-предложений.

\section{Таблица быстрой диагностики}

\begin{center}
\begin{tabular}{p{0.26\textwidth} p{0.28\textwidth} p{0.30\textwidth}}
\toprule
\textbf{Наблюдаемый симптом} & \textbf{Вероятный класс поломки} & \textbf{Первая проверка} \\
\midrule
ответы слишком общие & hub collapse & visits vs degree, top bridge nodes \\
контрпример лежит в графе, но не всплыл & bottleneck starvation & contradiction hit rate, conductance \\
много красивых выводов, мало объяснений & provenance erosion & provenance coverage, support-skeleton audit \\
часто уверена, потом резко переобувается & premature certainty & distribution of $\conf$, class priors \\
мрачные циклические ответы & mood lock-in & autocorrelation of mood, valence skew \\
контекст раздут и дорог & context bloat & evidence/glue ratio, token budget usage \\
\bottomrule
\end{tabular}
\end{center}

\section{Главный вывод}

\begin{principle}[Failure mode важнее красивой общей метафоры]
Пока у памяти нет таблицы собственных поломок, она остаётся литературным образом.
Как только у неё появляется каталог типичных сбоев, наблюдаемые симптомы и
предохранители, она становится инженерным объектом.
\end{principle}
